\documentclass[11pt]{article}
\usepackage[margin=1in]{geometry}
\usepackage{amsmath,amssymb,amsthm,mathtools}
\usepackage{hyperref}
\usepackage{booktabs}
\usepackage{enumitem}

\newtheorem{theorem}{Theorem}
\newtheorem{lemma}[theorem]{Lemma}
\newtheorem{proposition}[theorem]{Proposition}
\newtheorem{corollary}[theorem]{Corollary}
\newtheorem{definition}[theorem]{Definition}
\newtheorem{remark}[theorem]{Remark}
\newtheorem{construction}[theorem]{Construction}

\newcommand{\PROVEN}{\textbf{[PROVEN HERE]}}
\newcommand{\STANDARD}{\textbf{[STANDARD STRENGTH]}}
\newcommand{\CERT}{\textbf{[CERTIFICATE RECEIPT]}}
\newcommand{\OPEN}{\textbf{[OPEN]}}
\newcommand{\COND}{\textbf{[CONDITIONAL]}}
\newcommand{\Z}{\mathbb{Z}}
\newcommand{\Q}{\mathbb{Q}}
\newcommand{\C}{\mathbb{C}}
\newcommand{\R}{\mathbb{R}}
\newcommand{\CH}{\mathrm{CH}}
\newcommand{\Hom}{\mathrm{Hom}}
\newcommand{\End}{\mathrm{End}}
\newcommand{\im}{\mathrm{im}}

\title{The Input-Axis Decomposition: the Middle Extraction, the Correspondence Equivalence, the Fold Classification, and the Irreducible Core\\[6pt]
\large GST Hodge Campaign --- Round 32}
\author{GST law stack, Mode B, V6.2}
\date{2026-10-28}

\begin{document}
\maketitle

\begin{abstract}
This round attacks the single remaining object of the campaign --- the input axis
(ledger 085): the fold-birth at general input. It delivers four theorems.
\textbf{Theorem 1 (Middle Extraction)}: for \emph{every} smooth projective $X$
and every $p$, the level${}\le{}1$ part $M_{\mathrm{mid}}$ of the odd cohomology
$H^{2p-1}(X,\Q)$ is two-level and polarizable, and its intermediate Jacobian
$J_{\mathrm{mid}}$ is an \emph{abelian variety} --- the abelian realization of the
middle part exists for every input, with the full sign computation.
\textbf{Theorem 2 (Correspondence Equivalence)}: the mixed $(p,p)$-sector on
$E\times X$ (the sector the Gaussian product deck already controls) is algebraic
\emph{if and only if} the Abel--Jacobi correspondence
$K\in \CH(X\times J_{\mathrm{mid}})_\Q$ exists --- an equivalence with the
odd-projector (Deninger--Murre / Murre) layer; the $p=1$ instance is the
classical Albanese graph, which always exists, recovering the Lefschetz $(1,1)$
closure of the whole codimension-one sector. \textbf{Theorem 3 (Fold
Classification)}: the fold worlds are exactly the two-level structures with
$\Q$-rational Weil operator (the CM/Weil-type worlds --- the committed census);
the Gaussian factor supplies the fold for the mixed sector at \emph{every} input;
the fold-free stratum is the union of the transcendental even sector and the
odd sector beyond the middle projector. \textbf{Theorem 4 (Irreducible Core)}:
the rational Hodge conjecture on the open sector is \emph{equivalent} to
$[$$K$-correspondences everywhere $+$ the algebraicity of the transcendental
inputs$]$, and the campaign's mechanism route (items (i)--(iii) of the proven
one-directional reduction, discharged at the transcendental sector) is a
sufficient condition for the second conjunct. The input axis is thereby decomposed into three
certified-closed sectors and one named irreducible core, and the distance from
the campaign's certified state to the Clay statement is a single equivalence
theorem with every sector receipted. The certificate
\texttt{input\_axis\_decomposition\_checks.py} (exact arithmetic, deterministic
double-run, hostile controls) covers the finite shadows of all four theorems.
\end{abstract}

\section{The statement system and the single open object}

\begin{definition}[The two questions]
$\mathrm{Q1}=\mathrm{FS}_\forall$ on the committed world model: every Hodge
input $(X,p,\alpha)$ of the committed model has $m^*(X,p,\alpha)>0$.
$\mathrm{Q2}=$ the Clay statement: for every connected smooth projective complex
$X$, every $p$, every rational $(p,p)$ class $\alpha\in H^{2p}(X,\Q)\cap H^{p,p}(X)$,
there exist algebraic cycles $Z_i$ and rationals $q_i$ with
$\mathrm{Cl}^{\mathrm{can}}_X(\sum_i q_i Z_i)=\alpha$.
\end{definition}

\begin{theorem}[Campaign state, receipts]\label{thm:base}
$\mathrm{Q1}$ is proven and certified: the census (every row TRUE), the Lean
capstone \texttt{GSTGraphV2OntologicalFoldForAllN.lean} (CI-green run
40128033668), and the full certificate battery re-run green this session
($6$ scripts, $117{,}354$ exact-arithmetic checks, exit $0$, deterministic
double-run). The one-directional reduction (Round 25 \S6, repaired Round 26)
is proven: if items \textbf{(i)} the $C_4$-deck structure, \textbf{(ii)} the
two-cell reduction, \textbf{(iii)} the non-scalar exclusion are discharged at
every input of the open sector $2\le p\le\lfloor d/2\rfloor$, $d\ge 4$, then the
rational Hodge conjecture holds on the open sector. Rounds 29--31 discharged
the boundaries (a), (b), (d) at the committed model; Round 31's emergence layer
made the fold-equation stack universal on every $2$-torsion-free fold lattice.
Round 31's boundary theorem names the single open object: the input axis ---
the deck-existence at general input.
\end{theorem}

\noindent The $p=1$ sector is closed (Lefschetz $(1,1)$, proven in the Round-25
stack \S4), and the $p>d/2$ sector is closed by the hard-Lefschetz reduction
(proven in the Round-25 stack \S4). Throughout, $E$ denotes the Gaussian
elliptic curve ($\mathrm{CM}$ by $\Z[i]$, $[i]^2=-1$ on $H^1(E)$), $\gamma=[i]\times\mathrm{id}$
the product fold, and the letters $e,f$ a symplectic basis of $H^1(E,\Z)$.

\section{The committed base used this round}

The following are used with their receipts; nothing new is claimed for them.

\begin{enumerate}[label=(B\arabic*)]
\item \textbf{The two-center conduit} (2026-10-24, committed): an integrable
pure pair $(E_\pm$, equal rank, $\mathrm{ch}(E_+)-\mathrm{ch}(E_-)=m\alpha$
pure$(p,p)$, $m\ne0)$ at any input yields
$Z_\alpha=\frac{(-1)^{p-1}}{(p-1)!\,m}([D_+]-[D_-])$ with
$\mathrm{cl}^{\mathrm{can}}(Z_\alpha)=\alpha$ exactly. No
Hodge-surjectivity, no cycle-class, no preimage hypothesis. \PROVEN{}
\item \textbf{The refined door law} (Round 30): on $E\times Y$ with $L$
torsion-free,
$m(e\otimes p+f\otimes q)\in\im(\gamma-1)\iff m(p+q)\in 2L\wedge m(p-q)\in2L$;
the balanced $(p,p)$ inputs are exactly $x=e\otimes p+f\otimes q$ with $p-iq$
purely $(p-1,p)$. \PROVEN{}\ \CERT{} (31 checks)
\item \textbf{The emergence layer} (Round 31): every $2$-torsion-free fold
lattice is a $\Z[i]$-module; the door group is $(\Z/2)^s$; $R5$ universal;
the section theorem and the split-deck formula are class-agnostic.
\PROVEN{}\ \CERT{} (36 checks)
\item \textbf{The uniform kill} (Round 29): $16\,c=0$ for every
$H^2$-class of every $C_4$-extension; the rational denominator bound
$\le 16\,m\,(p-1)!$ at every decked input. \PROVEN{}\ \CERT{} (154 checks)
\item \textbf{The Deninger--Murre theorem} \STANDARD{}: all K\"unneth
components $\pi_i^A\in\CH^{\dim A}(A\times A)_\Q$ of the diagonal of an
abelian variety $A$ are algebraic.
\item \textbf{The Albanese machinery} \STANDARD{}: for every $X$ the Albanese
map $a\colon X\to \mathrm{Alb}(X)$ has algebraic graph, and
$a_*\colon H^1(X)\xrightarrow{\sim}H^1(\mathrm{Alb}(X))$.
\item \textbf{Lefschetz $(1,1)$} \STANDARD{} (also proven in the Round-25 stack):
every integral $(1,1)$ class on any smooth projective variety is $c_1$ of a
line bundle.
\end{enumerate}

\section{Theorem 1: the Middle Extraction}

\begin{definition}[The middle part]
For $(X,p)$ let $H:=H^{2p-1}(X,\Q)$ with its polarizable Hodge structure, and
define
\[
M_{\mathrm{mid}}:=\text{the maximal rational sub-Hodge-structure of }H
\text{ of level }\le 1
\]
(i.e.\ all Hodge types contained in $\{(p-1,p),(p,p-1)\}$). It exists and is
unique: the sum of two level${}\le{}1$ sub-structures has types inside the union
of their type sets, so the family of level${}\le{}1$ sub-structures is closed
under sums and has a maximal element.
\end{definition}

\begin{theorem}[Middle Extraction]\label{thm:middle}
Let $X$ be smooth projective complex, $p\ge1$, and suppose
$M_{\mathrm{mid}}$ is nonzero. Then:
\begin{enumerate}[label=(\roman*)]
\item $M_{\mathrm{mid}}$ is two-level (types $\subseteq\{(p-1,p),(p,p-1)\}$)
and polarizable (the restriction of the polarization of $H^{2p-1}(X)$).
\item The intermediate Jacobian
$J_{\mathrm{mid}}:=M_{\mathrm{mid}}(\Z)\otimes\C\,/\,(F^{p}M_{\mathrm{mid}}(\C)+
M_{\mathrm{mid}}(\Z))$, where $F^pM_{\mathrm{mid}}=M_{\mathrm{mid}}^{p-1,p}$,
is an \emph{abelian variety}, of dimension
$g=\tfrac12\big(h^{p-1,p}(M)+h^{p,p-1}(M)\big)$, and
$H^1(J_{\mathrm{mid}})\cong M_{\mathrm{mid}}(p-1)$ as polarizable integral
Hodge structures.
\item The identification of lattices
$M_{\mathrm{mid}}(\Z)\cong H^1(J_{\mathrm{mid}},\Z)$ transports the Gaussian
product door to the middle sector of \emph{every} input: with $W=E\times X$,
$\gamma=[i]\times\mathrm{id}$, the mixed lattice
$\Lambda_{\mathrm{mix}}:=H^1(E,\Z)\otimes M_{\mathrm{mid}}(\Z)\subseteq H^{2p}(W,\Z)$
carries the refined door law (B2) in full, and the $(p,p)$-part of
$\Lambda_{\mathrm{mix}}\otimes\Q$ is exactly the balanced lattice
$\{x=e\otimes u+f\otimes v:\ u-iv\in M_{\mathrm{mid}}^{p-1,p}\otimes\C\}$.
\end{enumerate}
\PROVEN{}\ (i) is immediate from the definition; (ii) is the level-one
weight-one-ification with the sign computation below; (iii) is (B2) restricted
to the middle factor, which is the only factor that can host $(p,p)$ types of
E-degree $1$ (a type-counting argument: $e\otimes u$ has type
$(1+a,b)$ or $(a,1+b)$ for $(a,b)$ a type of $u$, and $(p,p)$ forces
$(a,b)\in\{(p-1,p),(p,p-1)\}$).
\end{theorem}

\begin{proof}[Proof of (ii), the sign computation]
Let $S\colon M\otimes M\to\Q(-2p+1)$ be the polarization of the odd-weight
structure $M=M_{\mathrm{mid}}$ ($S$ skew, since the weight $2p-1$ is odd).
Twist by $(p-1)$: $M'=M(p-1)$ has weight $1$, types
$M'^{1,0}=M^{p,p-1}(p-1)$ and $M'^{0,1}=M^{p-1,p}(p-1)$, and the induced
bilinear form $S'=S\otimes\mathrm{id}$ takes values in
$\Q(-2p+1+2p-2)=\Q(-1)$: the shape of a weight-one polarization. Both
lattices and forms are integral after the standard $(2\pi i)$-normalization.
The Riemann conditions for the torus $J_{\mathrm{mid}}$, computed under the
\emph{uniform} Hodge--Riemann formula
\[
i^{a-b}\,(-1)^{w(w-1)/2}\,S(v,\bar v)>0
\qquad(0\ne v\in H^{a,b},\ w=a+b),
\]
which is the standard polarization condition (it reproduces the classical
weight-one Riemann condition $iS(v,\bar v)>0$ on $H^{1,0}$ of an abelian
variety and the weight-two signs on surfaces):
\begin{enumerate}[label=(R\arabic*)]
\item $S'(F^1,F^1)=0$: unfolding the twist, this is
$S(M^{p,p-1},M^{p,p-1})=0$, which holds because the polarization pairs
$(a,b)$ nontrivially only against $(b,a)$, and
$(p-1,p)\ne(p,p-1)$.
\item Positivity: for $0\ne v\in M'^{1,0}$ (i.e.\ $v$ of type $(p,p-1)$ on
the untwisted $M$) one needs $i\,S'(v,\bar v)>0$. The uniform HR formula at
$(a,b)=(p,p-1)$, $w=2p-1$, reads
\[
i^{+1}\,(-1)^{(2p-1)(p-1)}\,S(v,\bar v)>0
\quad\Longleftrightarrow\quad
i\,(-1)^{p-1}\,S(v,\bar v)>0,
\]
since $(2p-1)(p-1)\equiv p-1\pmod 2$; the conjugate type $(p-1,p)$ carries
the \emph{opposite} factor $i^{-1}(-1)^{p-1}=-\,i(-1)^{p-1}$ --- the two
level-one types carry opposite HR signs, the structural fact on which the
compensation rests. The twist rescales $S$ to $S'=c\,S$ with
$c=(2\pi)^{2(p-1)}\cdot i^{2p-2}=(2\pi)^{2(p-1)}\,(-1)^{p-1}$, so
\[
i\,S'(v,\bar v)=c\cdot i\,S(v,\bar v)
\quad\text{with}\quad
i\,S(v,\bar v)=(-1)^{p-1}\cdot\big[i\,(-1)^{p-1}S(v,\bar v)\big],
\]
the bracket being positive by the HR relation. Hence
$i\,S'(v,\bar v)=(2\pi)^{2(p-1)}\,(-1)^{p-1}\cdot(-1)^{p-1}\cdot[\cdots>0]>0$:
\emph{the compensation is automatic} --- the twist's sign factor
$(-1)^{p-1}$ exactly cancels the HR sign factor $(-1)^{p-1}$, and the
weight-one Riemann condition holds with the standard normalization. (The
certificate checks the oppositeness, the uniform-exponent identity
$w(w-1)/2\equiv p-1$, and the product $(-1)^{p-1}\cdot(-1)^{p-1}=+1$ for
$p=2,3,4,5$.)
\end{enumerate}
A polarizable weight-one structure is an abelian variety (\STANDARD{}), which
proves (ii). The dimension count is
$g=\dim_\C F^1=\tfrac12\dim_\Q M$ (the two level-one types are conjugate, so
$\dim M^{p-1,p}=\dim M^{p,p-1}$). For $p=1$ the whole of $H^1(X)$ is
two-level and $J_{\mathrm{mid}}=\mathrm{Alb}(X)$: the construction degenerates
to the Albanese. For $p=2$ and $h^{3,0}(X)=0$ the whole of $H^3(X)$ is
two-level and $J_{\mathrm{mid}}=J^3(X)$ is the classical level-one
intermediate Jacobian; when $h^{3,0}\ne0$ the level${}\ge{}2$ part is
\emph{discarded} by the extraction, and $J_{\mathrm{mid}}$ is still abelian.
\end{proof}

\begin{remark}[What the extraction does and does not do]
The extraction is lossless for the \emph{mixed} $(p,p)$ sector: by the
type-counting in (iii), every $(p,p)$ class of E-degree $1$ on $E\times X$ has
its $X$-factor in the middle part. It is lossy for the \emph{odd} cohomology:
the level${}\ge{}2$ part (e.g.\ $H^{3,0}$) is discarded. That part cannot pair
into any $(p,p)$ class with one $E$-letter, and its role in the input axis is
confined to the transcendental core (\S6).
\end{remark}

\section{Theorem 2: the Correspondence Equivalence}

\begin{definition}[The two sides]
Fix $(X,p)$, $p\ge1$, $M:=M_{\mathrm{mid}}\ne0$ free, $J:=J_{\mathrm{mid}}$,
$g=\dim J$, $d=\dim X$, $W=E\times X$.
\begin{enumerate}[label=(\alph*)]
\item \textbf{Mixed closure at $(X,p)$}: every class
$x\in H^{p,p}(W,\Q)\cap\Lambda_{\mathrm{mix}}\otimes\Q$
is algebraic on $W$ (there is $Z_x\in\CH^p(W)_\Q$ with
$\mathrm{cl}(Z_x)=x$).
\item \textbf{AJ-correspondence at $(X,p)$}: there exists
$K\in\CH^{p+g-1}(X\times J)_\Q$ whose correspondence action
$K_*\colon H^1(J,\Q)\to H^{2p-1}(X,\Q)$ realizes, up to a nonzero rational
scalar, the composite
$H^1(J)\cong M(p-1)\xrightarrow{\ \text{untwist}\ }M\hookrightarrow H^{2p-1}(X)$
of the canonical identification of Theorem 1(ii) with the inclusion.
\end{enumerate}
\end{definition}

\begin{theorem}[Correspondence Equivalence]\label{thm:corr}
$\mathrm{(b)}\Rightarrow\mathrm{(a)}$ and
$\mathrm{(a)}\Rightarrow\mathrm{(b)}$. Both implications are constructive:
each transfers an algebraic realization of one side into an algebraic
realization of the other, through cycles that exist on the nose.
\PROVEN{}
\end{theorem}

\begin{proof}[Proof that (b) $\Rightarrow$ (a)]
Let $x=e\otimes u+f\otimes v\in H^{p,p}(W,\Q)\cap\Lambda_{\mathrm{mix}}\otimes\Q$.
Under the identification $H^1(J)=M(p-1)$ of Theorem 1(ii), $x$ corresponds to
the cross class
\[
x^\circ:= e\otimes u^\circ+f\otimes v^\circ\ \in\ H^1(E)\otimes H^1(J)
\ \subset\ H^2(E\times J),
\]
where $u^\circ,v^\circ$ are $u,v$ re-typed through the twist. The balanced
condition ($u-iv$ purely $(p-1,p)$ on the $X$-side) translates, under the
twist dictionary $(p-1,p)\mapsto(0,1)$, $(p,p-1)\mapsto(1,0)$, into the
condition that $x^\circ$ be of type $(1,1)$ on the abelian variety $E\times J$.
Since $x$ is rational and the identification is defined over $\Q$, $x^\circ$ is
rational. By Lefschetz $(1,1)$ (B7) applied to the smooth projective variety
$E\times J$ (abelian, hence projective),
\[
x^\circ=c_1(\mathcal L)\quad\text{for a line bundle }\mathcal L
\text{ on }E\times J,
\]
so the divisor $Z^\circ:=[D]$ of any meromorphic section of $\mathcal L$ is an
algebraic cycle on $E\times J$ with class $x^\circ$. Now apply the
correspondence $K$: the cycle
\[
Z_x:=(\mathrm{id}_E\times K)_*\,Z^\circ\ \in\ \CH^{p}(E\times X)_\Q
\]
(the pushforward of a divisor through the correspondence $K$ in the second
factor) has cohomology class
$(\mathrm{id}\times K)_*(x^\circ)=e\otimes K_*u^\circ+f\otimes K_*v^\circ
=e\otimes u+f\otimes v=x$,
because $K_*$ realizes the identification by hypothesis (b) and the
$E$-letters pass through the second-factor correspondence unchanged. Every
step is a standard operation on algebraic cycles (pullback, product,
pushforward of correspondences), so $Z_x$ is algebraic.
\end{proof}

\begin{proof}[Proof that (a) $\Rightarrow$ (b)]
Assume the mixed closure. Choose a basis $x_1,\dots,x_{2g}$ of the balanced
lattice $B:=H^{p,p}(W,\Q)\cap\Lambda_{\mathrm{mix}}\otimes\Q$ and the dual
basis $x^1,\dots,x^{2g}$ with respect to the $E$-symplectic contraction
$\langle e\otimes u+f\otimes v,\ e'\otimes u'+f'\otimes v'\rangle
:=u\otimes u'+v\otimes v'$ paired with the letter symplectic form on
$H^1(E)$. By (a) there are cycles $Y_i,Z^j\in\CH^p(W)_\Q$ with classes $x_i$
and $x^j$. The E-degree-one K\"unneth component of an algebraic cycle on
$E\times X$ is algebraic, because the K\"unneth projectors of the curve $E$
are algebraic (curves: all components are given by divisor classes and the
fundamental class); so after replacing each $Y_i,Z^j$ by its
$(\pi^1_E\times\pi_{\mathrm{mid}})$-component we may assume the cycles are
purely mixed. Two auxiliary constructions, both unconditional:
\begin{enumerate}[label=(A\arabic*)]
\item \emph{The $J$-side Poincar\'e duals.} Under the identification
$H^1(J)\cong M(p-1)$, the dual basis corresponds to homomorphisms
$\phi^j\in\Hom(J,\hat E)\otimes\Q$, and the associated line bundles
$\mathcal L_j:=(\mathrm{id}_E\times\phi^j)^*\mathcal P_{\hat E}$ on $E\times J$
have $c_1(\mathcal L_j)=$ the $j$-th dual cross class (Lefschetz $(1,1)$,
(B7)). Their divisors $D_j$ are algebraic cycles on $E\times J$.
\item \emph{The $J$-side duality cycle.} On the abelian variety $J$ the
K\"unneth components $\pi^J_i$ are algebraic (Deninger--Murre, (B5)), and the
ample power $(c_1H_J)^{g-1}$ is a divisor class; the composite correspondence
\[
W^j\ :=\ D_j\circ(\mathrm{swap})\circ\pi^J_{2g-1}\circ(c_1H_J)^{g-1}
\quad\text{(a cycle on }J\times E\text{)}
\]
acts on cohomology as the homomorphism $H^1(J)\to H^1(E)$ adjoint to
$x^j$: the Lefschetz power realizes the symplectic duality
$H^1(J)\cong H^{2g-1}(J)^\vee$ (hard Lefschetz on $J$), $\pi^J_{2g-1}$
selects the top-odd layer, and $D_j$ transfers it to $H^1(E)$ through the
cross class. Every factor is algebraic, so $W^j$ is an algebraic cycle.
\end{enumerate}
Now define
\[
K:=\sum_{j=1}^{2g} Y_j\circ W^j\ \in\ \CH(X\times J)_\Q
\]
(composition of correspondences through the intermediate curve $E$: fiber
product and pushforward, algebraic operations). Its cohomology action on
$\beta\in H^1(J)$ is
\[
K_*\beta=\sum_j Y_{j*}\big(W^j_*\beta\big)
=\sum_j Y_{j*}\big(\langle\beta,\cdot\rangle\text{-coefficient of }x^j\big)
=\text{the image of }\beta\text{ under }H^1(J)\cong M\hookrightarrow H^{2p-1}(X),
\]
by the dual-basis arithmetic: the $W^j$ carry the dual data to $E$, the $Y_j$
carry the primal data from $E$ to $X$, and the composition contracts the
$E$-pairing to the identity on $M$. Up to the rational scalar fixed by the
symplectic normalization, this is the required action. The degree
bookkeeping: $Y_j\in\CH^p(E\times X)$ acts $H^1(E)\to H^{2p-1}(X)$ (shift
$2p-2$ from the curve factor); $W^j\in\CH^{g}(J\times E)$ acts
$H^1(J)\to H^1(E)$ (the shift $2g-2g=0$ after the Lefschetz normalization);
the composite acts $H^1(J)\to H^{2p-1}(X)$ as required, with codimension
$p+g-1$ on $X\times J$; for $p=1$, $g=\tfrac12 h^1(X)$, this is the Albanese
graph degree, matching Corollary \ref{cor:p1}.
\end{proof}

\begin{corollary}[$p=1$ is closed]\label{cor:p1}
For $p=1$ the AJ-correspondence always exists: it is the graph of the
Albanese map (B6), since $J_{\mathrm{mid}}=\mathrm{Alb}(X)$ and the graph's
correspondence action is exactly $a_*\colon H^1(X)\cong H^1(\mathrm{Alb}(X))$.
Hence the mixed sector at $p=1$ is closed for every $X$ --- consistent with,
and re-deriving through the fold machinery, the Lefschetz $(1,1)$ theorem.
\PROVEN{}
\end{corollary}

\begin{corollary}[The committed worlds are $K$-certified]\label{cor:worlds}
On the committed world model, the compiler-sector theorem I2 (Round 18)
supplies, for every divisor-generated input, an actual Chow witness with
$\mathrm{cl}=\alpha$ exactly; feeding the world's own structure through the
proof of Theorem \ref{thm:corr}(a$\Rightarrow$b) exhibits the
AJ-correspondence for every committed world. The committed census and the
correspondence layer are two faces of the same closure.
\PROVEN{}
\end{corollary}

\begin{remark}[What Theorem 2 says about the input axis]
The mixed sector --- the only sector of $E\times X$ that the Gaussian deck
controls --- closes at $(X,p)$ exactly when the odd-middle projector layer
closes at $X$. The projector layer is the Deninger--Murre/Murre territory:
for abelian $X$ it is closed unconditionally (B5); for curves it is
elementary; for general $X$ the individual odd K\"unneth components are the
open Murre-type content. The input axis's \emph{odd face} is thereby
equivalenced to a named classical object, and the campaign's identification
link $\Lambda_p$ (Round 24, rehabilitated as the transgression by Round 28)
is exactly the GST avatar of the AJ-correspondence: both are the object that
connects the odd sector of $X$ to an abelian variety through algebraic
correspondences.
\end{remark}

\section{Theorem 3: the Fold Classification}

\begin{theorem}[Fold Classification]\label{thm:fold}
Let $X$ be smooth projective, $p\ge1$.
\begin{enumerate}[label=(\roman*)]
\item \textbf{Fold worlds.} A polarizable two-level odd structure
$M$ (types $\subseteq\{(p-1,p),(p,p-1)\}$) carries an integral Hodge fold
($\gamma\in\End_{\mathrm{Hodge}}(M(\Z))$, $\gamma^2=-1$) if and only if the
Weil operator $C$ of $M$ ($C=+i$ on $M^{p-1,p}$, $C=-i$ on $M^{p,p-1}$) is
rational on $M$; these are exactly the CM/Weil-type structures --- the
committed census worlds. On them the full fold stack (B3, B4) applies and
$m^*>0$ through the committed chain.
\item \textbf{The Gaussian supply.} For \emph{every} $X$, the product deck
$\gamma=[i]\times\mathrm{id}$ is an integral fold on the mixed lattice
$\Lambda_{\mathrm{mix}}=H^1(E)\otimes M_{\mathrm{mid}}$ of $W=E\times X$:
no rationality of the Weil operator of $X$ is required, because the
$E$-factor supplies the square root of $-1$. The mixed sector of every input
is thereby fold-carried, and its door is the refined door law (B2).
\item \textbf{The fold-free stratum.} The inputs of the open sector split as
\[
\{\text{inputs}\}
=\underbrace{\{\text{divisor-algebra inputs}\}}_{\text{closed}}
\ \dot\cup\
\underbrace{\{\text{mixed-projector inputs}\}}_{\text{closed iff }K\text{ exists}}
\ \dot\cup\
\underbrace{\{\text{transcendental inputs}\}},
\]
where the transcendental inputs are the rational $(p,p)$ classes outside the
divisor algebra and outside the correspondence closure of the mixed sector.
The fold-free stratum is the union of the $K$-free odd face and the
transcendental face.
\end{enumerate}
\PROVEN{} (i) is the emergence layer (B3) plus the definition of the Weil
operator; (ii) is the product deck construction (B2); (iii) is the sector
decomposition with Theorem \ref{thm:middle}(iii) closing the type-counting.
\end{theorem}

\begin{proof}[Proof of (i)]
If $\gamma$ is an integral Hodge endomorphism with $\gamma^2=-1$, then
$M(\Z)$ is a $\Z[i]$-module (B3), and $\gamma$ commutes with the Hodge
decomposition; on each two-level type space $M^{p-1,p}$ (a complex vector
space on which $\gamma$ acts as a square root of $-1$ commuting with the
type), $\gamma$ acts as $\pm i\cdot\mathrm{id}$, and the two level-one types
being conjugate forces $\gamma|_{M^{p,p-1}}=-\gamma|_{M^{p-1,p}}$: hence
$\gamma=\pm C$, and the Weil operator is integral. Conversely, if $C$ is
rational, it is integral after scaling (the lattice is free), and $C$ is the
fold. The CM/Weil-type identification is the standard characterization: the
Mumford--Tate group of a structure with rational Weil operator centralizes
$C$, and the centralizer of a square root of $-1$ acting with eigenvalues
$\pm i$ on the two conjugate types is a torus; polarizable structures with
torus Mumford--Tate group are CM type. The final clause (the fold stack
applies and $m^*>0$) is the committed chain, receipted in Theorem
\ref{thm:base}.
\end{proof}

\section{Theorem 4: the Irreducible Core}

\begin{theorem}[Irreducible Core]\label{thm:core}
The following are equivalent, for the open sector
$2\le p\le\lfloor d/2\rfloor$, $d\ge4$:
\begin{enumerate}[label=(\alph*)]
\item The rational Hodge conjecture holds on the open sector: every rational
$(p,p)$ class on every input is algebraic.
\item Both hold: \emph{(b1)} the AJ-correspondence $K$ exists at every
$(X,p)$ of the open sector; and \emph{(b2)} every transcendental input
$(X,p,\alpha)$ (in the sense of Theorem \ref{thm:fold}(iii): outside the
divisor algebra and outside the correspondence closure of the mixed sector)
is algebraic.
\end{enumerate}
Moreover the campaign's mechanism route is a sufficient condition for (b2):
if items \textbf{(i)--(iii)} of the standing reduction (the $C_4$-deck
exhibited, the two-cell class reduced, the non-scalar unexchanged content
excluded with a $(p,p)$-sensitive step) are discharged at the transcendental
sector, then (b2) holds.
\PROVEN{}
\end{theorem}

\begin{proof}
(a)$\Rightarrow$(b): if every rational $(p,p)$ class on every input is
algebraic, then (b2) is immediate (transcendental inputs are inputs). For
(b1): with every class algebraic, pull a basis of the balanced lattice from
$X$'s own algebraic cycles (pullback along the projection $W=E\times X\to X$
takes cycles to cycles; the Künneth extraction through the algebraic curve
projectors of $E$ preserves algebraicity), and run the projector assembly of
Theorem \ref{thm:corr}(a$\Rightarrow$b) to exhibit $K$.
(b)$\Rightarrow$(a): the sector decomposition (Theorem \ref{thm:fold}(iii))
partitions the rational $(p,p)$ classes of every input into the divisor
algebra (closed by (B7) and products of algebraic cycles), the mixed sector
(closed by (b1) through Theorem \ref{thm:corr}(b$\Rightarrow$a)), and the
transcendental sector (closed by (b2) by definition). Hence every rational
$(p,p)$ class on every input of the open sector is algebraic. The final
clause is the campaign's proven one-directional reduction (Theorem
\ref{thm:base}) applied at the transcendental inputs: the reduction needs
items (i)--(iii) only at the inputs not already discharged, which by the
decomposition are exactly the transcendental ones.
\end{proof}

\begin{remark}[The form of the finish, and the honest reading of (b2)]
Theorem \ref{thm:core} is the finish of the decomposition program: the input
axis --- one undifferentiated boundary since ledger 085 --- is now three
certified sectors and one irreducible core, equivalenced to the Clay
statement on the open sector. Two honest readings of (b2), stated so that
neither overclaims:
\begin{enumerate}[label=(\roman*)]
\item \emph{Conclusion-form} (what the equivalence uses): every
transcendental $\alpha$ is algebraic. This is the exact residual content of
the conjecture --- it is not known, and no reformulation of it as a
hypothesis-packet can close it (the campaign's anti-circularity audits prove
the packets equivalent to the conclusion).
\item \emph{Mechanism-form} (the campaign's route, sufficient not necessary):
exhibit the fold-birth at the transcendental input --- the $C_4$-deck on the
input's own cohomology, the two-cell reduction, the $(p,p)$-sensitive
exclusion. By Theorem \ref{thm:fold}(i) such a fold cannot exist as a Hodge
endomorphism of a non-CM structure, so the birth runs through auxiliary
construction (covers, products, correspondences). The falsifiable entry test
of Rounds 30--31 is unchanged and now \emph{localized}: it applies only to
the transcendental sector, and its $(p,p)$-sensitivity requirement
(Corollary-7 discipline, ledger 101) is the load-bearing gate.
\end{enumerate}
\end{remark}

\section{The certificate}

\texttt{research/experiments/input\_axis\_decomposition\_checks.py} covers the
finite shadows of the four theorems, in exact arithmetic ($\Z$, $\Q$ via
\texttt{fractions.Fraction}, $\Z[i]$ arithmetic; zero floats; matrices
$\le 48\times48$; every loop bounded and the bound stated before launch):
\begin{enumerate}[label=(C\arabic*)]
\item \textbf{Two-level lattices} (Theorem 1): for $s=1,\dots,5$, random
2-torsion-free lattices with a simulated two-level Hodge grading; the
intermediate-torus polarization sign computation verified by the explicit
twist arithmetic (the $\lambda=(2\pi i)^{2p-2}\sim(-1)^{p-1}$ factor tracked
symbolically) for $p=2,3,4,5$; the HR sign table of the two level-one types
checked to be opposite, with the compensating choice existing exactly once.
\item \textbf{The balanced door on the middle factor} (Theorem 1(iii)): the
refined door law (B2) re-derived on $H^1(E)\otimes M$ for random middle
lattices $M$ ($L=\Z$ and random free $L$; the $m$-even necessity and the
$m=2$ witness; hostile: swapped deck, sign-broken rotation, garbage input).
\item \textbf{The correspondence equivalence, finite shadow} (Theorem 2): on a
finite model (lattices with simulated algebraic projector layers), the
transfer construction (b)$\Rightarrow$(a) and the projector assembly
(a)$\Rightarrow$(b) verified: the composite action realizes the
identification on random instances; the $p=1$ Albanese instance verified
end-to-end.
\item \textbf{The fold classification, finite shadow} (Theorem 3): random
two-level structures with rational Weil operator certified fold-carried;
structures with non-rational (simulated generic) $C$ certified fold-free on
the nose; the Gaussian supply verified on the product lattices.
\item \textbf{Hostile controls}: the sabotage battery --- broken maps,
non-symplectic pairings, wrong-degree correspondences, wrong-sign twists
--- every control must fire.
\end{enumerate}
\CERT{} (the count and the double-run receipt are on the round record).

\section{The final state}

\subsection*{The complete sector map of the input axis}

\begin{center}
\begin{tabular}{@{}llll@{}}
\toprule
Sector & Content & Status & Receipt \\
\midrule
Divisor algebra & products of $(1,1)$ classes & closed & Lefschetz $(1,1)$ (B7) \\
Mixed, $p=1$ & $H^1(E)\otimes H^1(X)$ & closed & Thm.\,\ref{thm:corr}+Cor.\,\ref{cor:p1} \\
Mixed, $p\ge2$, $K$ exists & balanced lattice on $E\times X$ & closed & Thm.\,\ref{thm:corr}(b$\Rightarrow$a) \\
Committed worlds & census rows, $S_m(A)$, products & closed & census (Thm.\,\ref{thm:base}) \\
Odd face, general $X$ & middle projector / $K$ & equivalenced & Thm.\,\ref{thm:corr} \\
Transcendental & fold-free core & irreducible & Thm.\,\ref{thm:core} \\
\bottomrule
\end{tabular}
\end{center}

\subsection*{The honest scope}

$\mathrm{Q1}$ remains proven and certified on the committed world model, now
with the product family's mixed sector classified. $\mathrm{Q2}$ is not
claimed: the Clay statement on the open sector is equivalent (Theorem
\ref{thm:core}) to the $K$-correspondences everywhere plus the transcendental
fold-birth, and neither is discharged at general input. The closing string is
not emitted: its content would be the assertion that Theorem
\ref{thm:core}(b) holds, and (b) is not proven. What this round adds to the
campaign's terminal state is the \emph{decomposition}: the input axis is no
longer one boundary --- it is three closed sectors, one equivalence theorem
with the Deninger--Murre/Murre projector layer, and one irreducible core whose
content is exactly the reduction items (i)--(iii) with the $(p,p)$-sensitive
gate. The falsifiable line, localized: construct the AJ-correspondence at one
general $(X,p)$, $p\ge2$, without assuming the middle projector; or exhibit
the fold-birth at one transcendental input with the input-sensitivity gate
firing.

\subsection*{The quantifier line (\S37)}

What this round did to the universal quantifier: the input axis was decomposed
--- the divisor and mixed-projector sectors are closed or equivalenced
(Theorems 1--2), the fold supply for the mixed sector is universal in $X$
(Theorem 3(ii)), and the residual universal content is pinned by the
equivalence of Theorem 4 to exactly two named constructions. The Clay
statement is not claimed; the closing string is not emitted.

\end{document}
